\documentclass[../../../../main.tex]{subfiles}
\begin{document}

\begin{flushright}
\footnotesize\textit{【自動文字起こし・要確認】}
\end{flushright}

\begin{proof}[解]
$x,y,z,w>0\ \cdots$①。文字$k$に対し，$k^{\frac1m}=K$とおく。
\[
(X+Y)^n+(Z+W)^n=\left\{(X^n+Z^n)^{\frac1n}+(Y^n+W^n)^{\frac1n}\right\}^n \quad\cdots\text{②}
\]
②がsuffix$(m,n)$で成立するので（$m,n\in\mathbb{N}$），$n=2$での成立が必要。
\[
2XY+2ZW=2\sqrt{(X^2+Z^2)(Y^2+W^2)} \quad\cdots\text{③}
\]
①から，③の両辺正で2乗しても良く，
\[
(XY+ZW)^2=(X^2+Z^2)(Y^2+W^2)
\]
\[
2XYZW=X^2W^2+Y^2Z^2
\]
\[
(XW-YZ)^2=0
\]
\[
XW=YZ \quad\cdots\text{④}
\]
が必要。逆に④が成立つ時，$\left(\dfrac{Z}{W}\right)=k\left(\dfrac{X}{Y}\right)$（$k\in\mathbb{R}_{>0}$）とおける。
\[
(\text{②の左辺})=(X+Y)^n+k^n(X+Y)^n=(1+k^n)(X+Y)^n
\]
\[
(\text{②の右辺})=\left[(1+k^n)^{\frac1n}X^{\frac1n\cdot n}+(1+k^n)^{\frac1n}Y^{\frac1n\cdot n}\right]^n=(1+k^n)(X+Y)^n
\]
とかり，たしかに②は成立して十分。

以上から
\[
XW=YZ
\]
つまり
\[
xw=yz
\]
\end{proof}

\newpage

\end{document}
